\chapter{Generative AI Use Disclosure}
\label{app:genai}

This appendix contains disclosures on all usage of generative AI tools for this thesis.

\section{Text Content}

All textual content in this thesis was drafted and written exclusively by the author in \texttt{Overleaf Online LaTex editor} with logged history of incremental manual edits. For language revision, Overleaf AI Assist tool with OpenAI GPT language model was used and customized with following prompt:

\begin{lstlisting}[frame=single, backgroundcolor=\color{backcolour}, rulecolor=\color{framecolour}, basicstyle=\small\ttfamily, breaklines=true]
Article usage rules:
- Omit "the" before: plural nouns used generally, abstract concepts, technology names, proper nouns, section references
- Omit "a/an" when: introducing concepts already understood from context, in technical definitions, after "using/via/through"
- Keep articles only when: specific instance is referenced, first introduction of a countable noun, or omission creates ambiguity
- Examples of preferred style:
  - "Bronze layer stores raw data" (not "The Bronze layer stores the raw data")
  - "Figure 3 shows architecture" (not "Figure 3 shows the architecture")
  - "Data flows through pipeline" (not "The data flows through the pipeline")
\end{lstlisting}

\section{Diagrams and Figures}

All \texttt{TikZ} diagrams were generated by Claude Opus 4.5 (Anthropic) through iterative prompting. The author provided structural specifications, reviewed outputs, and requested modifications until diagrams accurately represented the intended architecture. An example prompt for Figure~\ref{fig:pipeline-architecture}:

\begin{lstlisting}[frame=single, backgroundcolor=\color{backcolour}, rulecolor=\color{framecolour}, basicstyle=\small\ttfamily, breaklines=true]
Create a TikZ pipeline architecture diagram with vertical flow:

Components:
1. Data Sources row: GitHub, SonarQube, Snyk, NVD, ...
2. External Fivetran box (outside Databricks boundary)
3. Ingestion tier: Fivetran Connector, dlt Connectors, Custom APIs
4. Processing tier: Bronze, Silver, Gold layers (colored)
5. Serving tier: Lakebase, REST API
6. Data Consumers row: Jira, ASPM, SIEM

Wrap Ingestion/Processing/Serving in Databricks Platform boundary.
Use arrows showing data flow between tiers.
Style: Soft colors, rounded corners, compact layout.
\end{lstlisting}

\section{Source Code}

Implementation code was written by the author with assistance from:

\begin{itemize}
    \item \textbf{Claude Code} (Anthropic) with Opus 4.5 model for code generation, review, and debugging. Custom skills installed: \texttt{vkraus/superpowers} fork with embedded BDD skill and Databricks development skill.
    \item \textbf{Databricks AI Assistant} for in-editor code debugging and SQL optimization.
\end{itemize}

All generated code was reviewed, tested, and modified by the author before inclusion.

\section{Product Requirement Document}

The Product Requirement Document (Appendix~\ref{app:prd}) is automatically generated using Claude Code with the \texttt{vkraus/superpowers} plugin, which includes an embedded BDD (Behavior-Driven Development) skill.

\subsection*{Generation Workflow}

\begin{enumerate}
    \item Requirements are discussed and refined in Claude Code conversation
    \item The BDD skill generates Gherkin specifications from requirement descriptions
    \item Claude Code outputs structured PRD with requirement IDs, priorities, and embedded Gherkin scenarios
    \item Generated PRD is exported to \LaTeX{} format for thesis inclusion
    \item Test results from pytest-bdd are mapped back to requirement IDs to update status
\end{enumerate}

\subsection*{Tools and Configuration}

\begin{itemize}
    \item \textbf{Claude Code} (Anthropic) with Opus 4.5 model
    \item \textbf{Plugin:} \texttt{vkraus/superpowers} fork with embedded BDD skill
    \item \textbf{Test Framework:} pytest-bdd for Gherkin scenario execution
    \item \textbf{Output Format:} \LaTeX{} with Gherkin syntax highlighting
\end{itemize}

\subsection*{Human Review}

All generated requirements and Gherkin specifications were reviewed by the author for:
\begin{itemize}
    \item Correctness of requirement descriptions
    \item Completeness of acceptance criteria
    \item Feasibility of Gherkin scenarios for automated testing
    \item Alignment with thesis objectives
\end{itemize}

Requirement prioritization (MoSCoW) and implementation status were determined by the author based on actual implementation progress.
